\documentclass[a4paper,9pt]{extarticle}

%=========================================================
% PACKAGES
%=========================================================

\usepackage[
a4paper,
top=1.8cm,
bottom=1.5cm,
footskip=7mm,
left=1cm,
right=1cm,
headheight=14pt,
headsep=5mm
]{geometry}
\usepackage[T1]{fontenc}
\usepackage[utf8]{inputenc}
\usepackage[french]{babel}

\usepackage[sfdefault]{FiraSans}
\renewcommand{\familydefault}{\sfdefault}

\usepackage{multicol}
\setlength{\columnsep}{0.7cm}

\usepackage{amsmath,amssymb}

\usepackage{titlesec}
\usepackage{enumitem}

\usepackage{tcolorbox}
\tcbuselibrary{skins,breakable}

\usepackage{xcolor}

\usepackage{graphicx}
\usepackage{float}

\usepackage{tikz}

\usepackage{fancyhdr}
\usepackage[hidelinks]{hyperref}


%=========================================================
% PAGE SETTINGS
%=========================================================

\setlength{\parindent}{0pt}
\setlength{\parskip}{1.5pt}

\setlength{\abovedisplayskip}{4pt}
\setlength{\belowdisplayskip}{4pt}

\setlist[itemize]{
leftmargin=*,
itemsep=1pt,
topsep=1pt
}


%=========================================================
% COLORS
%=========================================================

\definecolor{Blue}{HTML}{E3F0FF}
\definecolor{Green}{HTML}{E3F7EA}
\definecolor{Pink}{HTML}{FFE5ED}
\definecolor{Orange}{HTML}{FFF0DC}
\definecolor{Purple}{HTML}{EEE7FF}
\definecolor{Yellow}{HTML}{FFFADB}
\definecolor{Teal}{HTML}{E0F5F3}


\definecolor{BlueTitle}{HTML}{285B8F}
\definecolor{GreenTitle}{HTML}{287052}
\definecolor{PinkTitle}{HTML}{A84063}
\definecolor{OrangeTitle}{HTML}{A96500}
\definecolor{PurpleTitle}{HTML}{6245A5}
\definecolor{TealTitle}{HTML}{1F6F68}


%=========================================================
% HEADER
%=========================================================

\pagestyle{fancy}

\fancyhf{}

\fancyhead[L]{\small\textbf{Modèles stochastiques}}
\fancyhead[C]{\small Convergence de suites de v.a.}
\fancyhead[R]{\small Week 03 --- 21-27/09/2026}

\fancyfoot[C]{\small\href{https://github.com/AlphaXT38G/notes_ba5}{github.com/AlphaXT38G/notes\_ba5}}
\fancyfoot[R]{\small \thepage}

\renewcommand{\headrulewidth}{0.4pt}



%=========================================================
% TITLES
%=========================================================

\titleformat{\section}
{\large\bfseries\color{BlueTitle}}
{}{0pt}{}

\titlespacing{\section}
{0pt}{6pt}{3pt}


% separator

\newcommand{\separator}{
\vspace{2pt}
\hrule
\vspace{3pt}
}


%=========================================================
% BOXES
%=========================================================

% Boxes are kept whole (unbreakable): inside multicols, a box split by tcolorbox
% leaves an empty column behind it.

% DEFINITION - blue
\newtcolorbox{definition}{
enhanced,
unbreakable,
colback=Blue,
colframe=BlueTitle,
boxrule=0.7pt,
arc=0mm,
title=\textbf{Definition},
fonttitle=\bfseries
}

% IMPORTANT - pink
\newtcolorbox{important}{
enhanced,
unbreakable,
colback=Pink,
colframe=PinkTitle,
boxrule=0.7pt,
arc=0mm,
title=\textbf{Key idea},
fonttitle=\bfseries
}

% FORMULAS - green
\newtcolorbox{formula}{
enhanced,
unbreakable,
colback=Green,
colframe=GreenTitle,
boxrule=0.7pt,
arc=0mm,
title=\textbf{Key formulas},
fonttitle=\bfseries
}

% WARNING - orange
\newtcolorbox{warning}{
enhanced,
unbreakable,
colback=Orange,
colframe=OrangeTitle,
boxrule=0.7pt,
arc=0mm,
title=\textbf{Warning},
fonttitle=\bfseries
}

% EXAMPLE - purple
\newtcolorbox{example}{
enhanced,
unbreakable,
colback=Purple,
colframe=PurpleTitle,
boxrule=0.7pt,
arc=0mm,
title=\textbf{Example},
fonttitle=\bfseries
}

% SUMMARY - orange
\newtcolorbox{summary}{
enhanced,
unbreakable,
colback=Yellow,
colframe=OrangeTitle,
boxrule=0.7pt,
arc=0mm,
title=\textbf{Summary},
fonttitle=\bfseries
}



% SIMPLE - teal: the same idea again, in the plainest words, with an everyday picture
\newtcolorbox{simple}{
enhanced,
unbreakable,
colback=Teal,
colframe=TealTitle,
boxrule=0.7pt,
arc=0mm,
title=\textbf{Expliqué simplement},
fonttitle=\bfseries
}



%=========================================================
% FIGURE BOX
%=========================================================

\newtcolorbox{figurebox}{
enhanced,
colback=white,
colframe=gray,
boxrule=0.5pt,
arc=0mm,
center
}



%=========================================================
% DOCUMENT
%=========================================================

\begin{document}

%\footnotesize   --> To reduce the size of the text

\thispagestyle{fancy}



\begin{multicols}{2}


%=========================================================
\section{Où on en est}


\begin{summary}

\begin{itemize}

\item Mardi : module 2, convergence \textbf{presque sûre} de suites de v.a.

\item Vendredi : module 2 (fin), les autres convergences.

\item Exercices M2 : 5, 6, 10, 11 (mardi) ; 9, 19, 20 (vendredi).

\end{itemize}

\end{summary}


\begin{definition}

Une suite de v.a. $\{X_n\}_{n\ge1} = \{X_1, X_2, \dots\}$. On étudie sa limite quand $n \to
\infty$, comme une suite de réels en analyse. \textbf{4 notions} de convergence vers une v.a.
$X$ : presque sûre, en distribution, en probabilité, en moyenne quadratique.

\end{definition}


\begin{simple}

Une suite de réels, c'est facile : elle s'approche d'un nombre ou pas. Une suite de v.a., c'est
plus flou, parce que chaque $X_n$ est au hasard. « S'approcher » peut vouloir dire plusieurs
choses : chaque tirage s'approche ? la chance d'être loin devient petite ? l'erreur moyenne
devient petite ? juste la forme de la loi s'approche ? D'où 4 définitions.

\end{simple}


%=========================================================
\section{Convergence presque sûre}


\begin{definition}

En fixant un résultat $\zeta$, $\{X_n(\zeta)\}$ est une \textbf{suite de réels}. Si elle
converge vers $X(\zeta)$ pour \textbf{tout} $\zeta$ : convergence \textbf{sûre} (partout). Si
les $\zeta$ où elle ne converge pas ont une probabilité nulle :

\[
P\big(\{\zeta \mid \lim_{n\to\infty}X_n(\zeta) = X(\zeta)\}\big) = 1
\]

$\{X_n\}$ converge \textbf{presque sûrement} (presque partout, avec probabilité 1) :
$X_n \xrightarrow{\text{p.s.}} X$ (ou $\xrightarrow{\text{a.s.}}$, $\xrightarrow{\text{w.p.1}}$).

\end{definition}


\begin{important}

Cette probabilité est \textbf{difficile à calculer directement}. On passe par des critères
sur $\sum_{n=1}^\infty P(|X_n - X| \ge \varepsilon)$, pour tout $\varepsilon > 0$ :

\begin{enumerate}[leftmargin=*,itemsep=1pt,topsep=1pt]

\item continuité des probabilités ;

\item lemmes de Borel-Cantelli, pour une suite d'évènements $\{A_n\}$ ;

\item on les applique aux évènements $A_n = \{|X_n - X| \ge \varepsilon\}$.

\end{enumerate}

\end{important}


\begin{simple}

On lance l'expérience \textbf{une fois} : ça donne toute une suite de nombres $X_1, X_2,
\dots$ (un chemin). Presque sûre veut dire : \textbf{presque tous les chemins} finissent par
arriver à la cible. Il peut exister des chemins bizarres qui n'y arrivent pas, mais ils ont
une chance zéro de se produire.

\end{simple}


%=========================================================
\section{Continuité des probabilités}


\begin{formula}

\textbf{Lemme 1 (i)} Suite \textbf{croissante} $A_1 \subseteq A_2 \subseteq \dots$, de limite
$A = \bigcup_{n=1}^\infty A_n$ :

\[
P(A) = \lim_{n\to\infty}P(A_n)
\]

\textbf{Lemme 1 (ii)} Suite \textbf{décroissante} $B_1 \supseteq B_2 \supseteq \dots$, de limite
$B = \bigcap_{n=1}^\infty B_n$ :

\[
P(B) = \lim_{n\to\infty}P(B_n)
\]

\end{formula}


\begin{important}

\textbf{Preuve de (i).} On écrit $A$ comme une union d'évènements \textbf{disjoints} : $A =
A_1 \cup (A_2\setminus A_1) \cup (A_3 \setminus A_2) \cup \dots$ L'axiome A3' donne

\[
P(A) = P(A_1) + \lim_{n\to\infty}\sum_{k=1}^{n-1}\big(P(A_{k+1}) - P(A_k)\big)
\]
\[
= \lim_{n\to\infty} P(A_n)
\]

(une somme télescopique). (ii) se démontre pareil avec les compléments.

\end{important}


\begin{simple}

Si des évènements grossissent petit à petit (comme une tache d'encre qui s'étale), la chance
de la tache finale est la limite des chances des taches intermédiaires. Pas de saut surprise
à l'infini.

\end{simple}


%=========================================================
\section{Lemmes de Borel-Cantelli}


\begin{definition}

$\{A_n \text{ i.o.}\}$ (\emph{infinitely often}) : l'évènement « une infinité des $A_n$ se
produisent ».

\[
\{A_n \text{ i.o.}\} = \bigcap_{n=1}^\infty \bigcup_{k\ge n} A_k
\]

Car $\zeta$ est dans une infinité de $A_k$ $\Leftrightarrow$ pour tout $n$, $\zeta$ est dans au
moins un des $A_n, A_{n+1}, \dots$

\end{definition}


\begin{formula}

\textbf{Premier lemme} (toute suite d'évènements) :

\[
\sum_{n=1}^\infty P(A_n) < \infty \;\Longrightarrow\; P(A_n \text{ i.o.}) = 0
\]

\textbf{Second lemme} (évènements \textbf{indépendants}) :

\[
\sum_{n=1}^\infty P(A_n) = \infty \;\Longrightarrow\; P(A_n \text{ i.o.}) = 1
\]

\end{formula}


\begin{important}

\textbf{Preuve du premier.} $\bigcup_{k\ge n}A_k$ décroît avec $n$, donc par le lemme 1 (ii) :

\[
P(A_n\text{ i.o.}) = \lim_{n\to\infty}P\Big(\bigcup_{k\ge n}A_k\Big) \le \lim_{n\to\infty}
\sum_{k\ge n}P(A_k) = 0
\]

car le reste d'une série convergente tend vers 0.

\vspace{3pt}

\textbf{Preuve du second.} Par le complément, l'indépendance et $1 - x \le e^{-x}$ :

\[
P\Big(\bigcap_{k\ge n}\bar A_k\Big) = \prod_{k\ge n}(1-P(A_k)) \le e^{-\sum_{k\ge n}P(A_k)} = 0
\]

donc $P(A_n\text{ i.o.}) = 1$.

\end{important}


\begin{simple}

On lance des fléchettes, une par jour, et $A_n$ = « je touche le centre le jour $n$ ». Si la
somme de toutes les chances est \textbf{finie} (les chances baissent vite, comme $1/n^2$), on
ne touchera le centre qu'un nombre fini de fois : un jour, c'est fini pour toujours. Si la somme
est \textbf{infinie} (elles baissent lentement, comme $1/n$) et que les jours sont
indépendants, on le touchera une infinité de fois, c'est garanti.

\end{simple}

\end{multicols}


\end{document}
